\section{Conclusion and Future Work}\label{sec:conclusion}

In US adults, a lung-cancer history concentrates in older people who have ever smoked and in people with COPD -- even after accounting for smoking and age. The detailed profile of the population shows where the \nCases cases are: three in four are aged 65 or older, four in five have ever smoked, and almost half have COPD. Regional pollution and smoking levels add little at the resolution available.

\subsection{Next steps: the modelling phase (planned, not implemented)}\label{sec:next}

\begin{enumerate}
  \item \textbf{Baselines first.} An ``always no'' model and the USPSTF-style rule (\file{screening_age_smoker}). On the test rows the rule flags 21.9\% of adults, finds 66\% of the cases, and 1.25\% of the people it flags have a lung-cancer history -- three times the base rate (Table~\ref{tab:screening}). Any real model has to beat both.
  \item \textbf{Models.} Class-weighted logistic regression on the \nFeatures features -- the reference categories are already dropped and documented, and no feature has a VIF above \maxVIFshort{} -- for interpretable adjusted odds ratios, then tree-based models such as random forest and gradient boosting. Train on the fixed 80/20 split and tune with stratified 5-fold cross-validation on training rows only.
  \item \textbf{Metrics.} Recall, precision, F1 and PR-AUC, with a threshold tuned for high recall, because missing a real case costs more than a false alarm; calibration checked on the untouched test set. No accuracy.
  \item \textbf{Ablation tests.} Retrain without the reverse-causation-prone features and without the regional features, to measure what each group really contributes; compare original features with PCA components for distance-based models.
  \item \textbf{Close the survey-design gap.} Extend the weighted, cluster-robust analysis (\file{PSTRAT}, \file{PPSU}) from the three key odds ratios to every association and to the model's evaluation.
  \item \textbf{Use more of NHIS.} The raw files contain variables we did not use: age at lung-cancer diagnosis (\file{LUNGAGETC_A}, to separate recent from old diagnoses), family history of cancer and current smoking intensity.
  \item \textbf{Fairness check.} Compare recall across sex, race/ethnicity and income groups before anyone relies on the model.
  \item \textbf{Reporting standard.} Write up the model with the TRIPOD+AI checklist~\cite{tripod} and a leakage information sheet~\cite{kapoor}.
\end{enumerate}

\begin{table}[htbp]
\centering\small
\caption{The USPSTF-style rule (\texttt{screening\_age\_smoker}: aged 50--80 and ever smoked) as a benchmark: recall = share of cases flagged; precision = share of flagged adults who are cases; lift = precision divided by the base rate. Source: \texttt{screening\_rule\_baseline.csv}.}\label{tab:screening}
\begin{tabular}{lrrrrrr}
\toprule
Rows & Adults & Flagged & \makecell[b]{Cases\\{}flagged} & Recall & Precision & Lift \\
\midrule
All adults & 149,928 & 32,373 (21.6\%) & 389 of 619 & 62.8\% & 1.20\% & 2.9 \\
Test rows & 29,986 & 6,579 (21.9\%) & 82 of 124 & 66.1\% & 1.25\% & 3.0 \\
\bottomrule
\end{tabular}

\end{table}
